\subsection{Confirmation and first-message authentication}
\label{sec:confirmation-traffic}

The same native rekey uses the three confirmation keys and directional
traffic material of Theorem~\ref{thm:joint-message-leaf}. We now retain
each confirmation key behind its actual MAC and verification interfaces
and derive a joint authentication consequence. Public queries to the
independent oracles $H$ and $h$ may be quantum. Protocol requests, MAC
inputs, verification requests and their replies are classical.

\paragraph{Native transcripts and application interface.}
Write $\mathcal H_L(x)=H(\operatorname{LP}(D\Vert L)
\Vert\operatorname{LP}(x))$, with $\operatorname{LP}$ the source's
eight-byte length prefix. For $d\in\{r,f,a\}$, let $C_d$ be the responder,
final or receipt confirmation key. The tag is
\begin{equation}
\begin{aligned}
 t_d&=F(C_d,\mathcal H_{L_d}(X_d)),\\
 (L_r,L_f,L_a)&=(\texttt{response-core},\texttt{final-core},
                         \texttt{receipt-core}).
\end{aligned}
\label{eq:native-confirmation-tag}
\end{equation}
The complete responder core contains the control tag, profile, context,
session, current and target epochs, sender role, control parent, digest of
the complete signed offer, and hybrid ciphertext. The final core contains
its corresponding control prefix, the offer and response wire digests,
and the old sending-chain length. The receipt additionally contains the
complete final-wire digest. Their respective source lengths are 1305,
225 and 257 bytes. Every MAC input in~\eqref{eq:native-confirmation-tag}
has 32 bytes.

The native verifier checks the complete prefix, body width and identity
signature before the relevant MAC verification. The public result follows
the actual success, failure, persistence and status transitions. The MAC
comparison may therefore give the reduction the expected tag internally;
the verification caller obtains the native result. Every honest tag
construction enters a ghost origin ledger at its first private
computation, before a subsequent write or release. Verification
computations have their own query ledger. Exact repetitions and matching
endpoint aliases use the same role-key identity and the same original
construction entry. A different signature encoding of an already logged
core retains that core's origin.

For the marked traffic allocation and direction, use the first-message
interface of Corollary~\ref{cor:joint-first-message}. The encryption nonce
is $0^{96}$, and the associated data is
\[
 D_M\Vert\texttt{aead}\Vert\mathit{header}_{93}
       \Vert\operatorname{BE16}(|ad|)\Vert ad.
\]
The header binds session, direction, epoch, index, message identifier and
plaintext length. The first private encryption is recorded by its exact
$(\mathit{AAD},\mathit{ciphertext})$ pair and original send identity.
Retries retrieve its cached result through Resume. Plaintext delivery
belongs to the private application. A fresh authenticated receive means
a successful AEAD check followed by installation of the corresponding
native inbox entry. The source's \texttt{PreviouslyConsumed} result is
a historical consumption status. An existing inbox entry retains its
original authenticated receive-intent digest and plaintext; a cached
return refers to that earlier receive.

Let $\mathsf{Bad}_{\rm conf}$ mean that, before the declared local stop,
a successful native confirmation verification accepts a complete core
with no earlier honest private tag construction under that role-key
identity. Let $\mathsf{Bad}_{\rm msg}$ mean that a fresh authenticated
receive under the selected first-message key installs a pair absent from
its honest encryption ledger. These definitions include failures and
unaccepted executions in the probability space.

\paragraph{Operational bounds.}
Fix a role order $r,f,a$. After exact-input caching, let $T_d$ bound the
distinct 32-byte messages evaluated under $C_d$ by both tag construction
and verification, and let $V_d$ bound verification attempts. Repeated
guesses count in $V_d$. All budgets are fixed before execution. Define
\begin{equation}
 \Delta_{\rm conf}=\min\left\{1,
       \sum_{d\in\{r,f,a\}}
          \varepsilon_{\rm Hrow}(T_d,q_d,32)\right\},
 \qquad V=V_r+V_f+V_a.
\label{eq:confirmation-row-cost}
\end{equation}
Here $q_d$ counts every ordinary $h$ query of that row's complete
distinguisher. In particular, it includes preparation, the target
application, all other real primitives and terminal testing. If $B_c$
bounds all such work outside these three confirmation MAC rows, the
conservative choices
\[
 q_d=B_c+4\sum_{e>d}T_e
\]
are valid: earlier rows are already ideal, later rows remain real, and a
32-byte HMAC evaluation uses at most four compression calls including
its two pads. The row lemma separately charges the replaced row's
private evaluations.

Let $\chi_H(Q)$ bound production of two distinct complete strings with
equal $H$ output by a $Q$-query algorithm. It counts the entire derived
experiment's ordinary and honest $H$ evaluations, including preparation
and the final collision test. For the random oracle with 256-bit output,
an admissible bound is
\[
 \chi_H(Q)=\min\{1,C_{\rm col}(Q+2)^3 2^{-256}\},
\]
with the universal constant of the quantum collision theorem
\cite[Theorem~7]{zhandryCollision2015}. The operational notation retains
that constant.

Finally, let $\varepsilon_{\rm INT}$ bound a quantum adversary's success
in the one-key, at-most-one-encryption ciphertext-integrity experiment
with nonce $0^{96}$, classical decryption and the same independent public
oracles. A successful decryption of a pair absent from the encryption
ledger wins. All repeated decryption attempts are included in its query
budget. The encryption ledger begins at the first private computation.

\begin{proposition}[Joint confirmation and message authentication]
\label{prop:confirmation-message-auth}
Use the controller of Theorem~\ref{thm:joint-message-leaf} and the above
native interfaces. Each $C_d$ is accessed through its tag and verification
service; the selected message key is accessed through its single original
encryption, classical decryption and exact-retry service. Then
\begin{equation}
\begin{split}
 \Pr_{G_R}[\mathsf{Bad}_{\rm conf}\lor\mathsf{Bad}_{\rm msg}]
 \le\min\{1,\;&\Delta_{\rm sched}+\Delta_{\rm conf}
       +V2^{-256}\\
       &+\chi_H(Q_H^{\rm conf})+\varepsilon_{\rm INT}\}.
\end{split}
\label{eq:confirmation-message-auth}
\end{equation}
The same controller retains the qualified confidentiality bound
$|\mathbb E_{G_R}Z|\le2\Delta_{\rm sched}
 +2\varepsilon_{\rm AEAD}$ of
Corollary~\ref{cor:joint-first-message}. Its confirmation and verification
operations are part of that complete continuation.
\end{proposition}
\begin{proof}
Apply Theorem~\ref{thm:joint-message-leaf} to the union event's terminal
bit. This costs $\Delta_{\rm sched}$ once and provides the joint
independent-key experiment $G_{\rm leaf}$. All transcript strings and
subsequent information strings are generated in the current world.

For confirmation, replace the three HMAC rows in the fixed role order by
independent private random functions $W_d$. The adaptive classical row
lemma applies to every tag and verification evaluation; its ordinary
query budgets include the remaining real rows. This costs
$\Delta_{\rm conf}$. Independent key samples may coincide numerically;
the real rows retain the shared compression oracle, and the row
comparisons account for those coincidences.

Consider a second execution that returns failure on every verification
whose digest has no earlier honest tag-construction entry under that
role. The coupling stops at the first internal MAC equality on such an
unissued digest, before subsequent persistence or a public response.
Up to this first hit, the two executions have
the same classical replies and induce the same transformation of the
adversary's quantum private state. In the second execution, the unissued
value $W_d(m)$ is independent of the available prefix: its verification
result is fixed to failure and its table value remains private. Each
of at most $V$ attempted tags equals this value with probability
$2^{-256}$. A union bound therefore bounds the first fresh-digest
success by $V2^{-256}$. It also covers repeated guesses for the same
digest. If a fresh complete core instead has a digest already used by
an honest construction under that role, the two different framed core
strings give a collision for $H$. The reduction records these classical
inputs and outputs the pair; its probability is bounded by
$\chi_H(Q_H^{\rm conf})$.

For message integrity in $G_{\rm leaf}$, sample every key except the
selected message key locally, preserving the actual confirmation, ACK,
successor-root and subsequent-chain computation. Forward the selected
original encryption and every selected-key decryption to the integrity
challenger. Exact retries use the first private encryption cache. A new
authenticated inbox installation without an encryption origin entails a
successful new decryption pair, so its probability is at most
$\varepsilon_{\rm INT}$. Existing inbox entries and historical consumed
statuses use their original metadata. The union bound in
$G_{\rm leaf}$ proves~\eqref{eq:confirmation-message-auth}.

For confidentiality, the AEAD reduction samples the confirmation keys
locally and executes their actual verification services. Thus
Corollary~\ref{cor:joint-first-message} applies with its unchanged
complete-continuation budget. For an independent application challenge
bit, both arguments may first fix that bit and then average.
\end{proof}

The authentication statement starts at the uniform epoch-root interface
$G_R$. The established root-initialization and complete-tuple comparisons
can precede it at their output-bit costs, for a controller realizing their
stated interfaces. A native authentication event additionally uses the
record and component comparisons for that event's own terminal bit.
The component costs $a_j^+,a_j^-$ in the confidentiality theorem are the
specified signed-score event distinguishers. This keeps the accounting
of a new authentication event explicit.
